\chapter{Timing and allowance accounting}\label{module:timing}
\Q{Evidence}{Current implementation boundary}
\path{Software/DemoCode/timing.py}. Exact file hashes are in the accompanying source manifest. Names, declarations and source locations below come from the captured source, not a runtime import.
\Q{Purpose}{Input-to-output reading flow}
Measure elapsed stage intervals and aggregate a request record. Compare those measurements with nominal motion allowances only after checking the stage boundary, host and completion status.


\Q{Evidence}{Selected current module settings}
These are declarations or expressions, not recovered environment values. Source comments do not establish a measured optimum. Larger glyph coordinate tables are designed data in the bundled source.
\begin{lstlisting}[style=plain]
ALLOWANCE = {'classify': lambda chars: 2.0 + 0.025 * chars, 'reproduce': lambda chars: 4.0 * chars, 'draw': lambda chars: 4.0 * chars}
ALLOWANCE_TEXT = {'classify': '2 s + 0.025 s per character', 'reproduce': '4 s per character', 'draw': '4 s per character'}
\end{lstlisting}
\section*{Stage: Timing and requirement accounting}
\section{Method background: Timing and requirement accounting}\label{method:timing}
\Q{Purpose}{The role in this system}
Measure request wall time and retain earlier generation work when drawing is requested later.


\Q{Justification}{The reason for this construction}
perf\_counter differences measure elapsed time; Carry adds external work while Add only labels work already inside the wall interval. This avoids double counting stages.


\Q{Alternatives}{What changes if another method is used}
CPU time excludes waiting; summing nested stages can double-count; theoretical path time ignores overhead. A benchmark on another host is not an embedded qualification.


\Q{Library}{Implementation speed and accuracy}
Python perf\_counter is already the timing primitive. Profilers can attribute cost more finely but add overhead. Multiprocessing stage wall time is not the sum of worker CPU times. \cite{pythontime,code2026}.


\section{Stage}\label{fn:timing:Stage}
\begin{lstlisting}[style=sig]
class Stage
\end{lstlisting}
\Q{Evidence}{Source and method link}
\path{Software/DemoCode/timing.py} line 47. Implementation evidence: \cite{code2026}. Related method, alternatives and library discussion: Section~\ref{method:timing}.


\Q{Purpose}{What this contributes}
Measure request wall time and retain earlier generation work when drawing is requested later.


\Q{Method}{Methods exposed by this current class}
\begin{lstlisting}[style=plain]
def __init__(self, run, name, chars=0)
def __enter__(self)
def __exit__(self, *exc)
\end{lstlisting}
\Q{Method}{Current implementation}
Measure request wall time and retain earlier generation work when drawing is requested later. The excerpt below is from the current definition; the complete function is in the source snapshot. Its constants and branches, rather than an older signature, govern this path.


\begin{lstlisting}[style=plain]
class Stage:
    """See source snapshot for the original docstring."""

    def __init__(self, run, name, chars=0):
        self.run = run
        self.name = name
        self.chars = chars
        self.t0 = None

    def __enter__(self):
        self.t0 = time.perf_counter()
        return self

    def __exit__(self, *exc):
        wall = time.perf_counter()-self.t0
        self.run.Record(self.name, wall, self.chars)
        return False
\end{lstlisting}
\Q{Justification}{Interpret the method rather than infer a benchmark}
perf\_counter differences measure elapsed time; Carry adds external work while Add only labels work already inside the wall interval. This avoids double counting stages. \cite{pythontime,code2026}.


\section{Run}\label{fn:timing:Run}
\begin{lstlisting}[style=sig]
class Run
\end{lstlisting}
\Q{Evidence}{Source and method link}
\path{Software/DemoCode/timing.py} line 66. Implementation evidence: \cite{code2026}. Related method, alternatives and library discussion: Section~\ref{method:timing}.


\Q{Purpose}{What this contributes}
Measure request wall time and retain earlier generation work when drawing is requested later.


\Q{Method}{Methods exposed by this current class}
\begin{lstlisting}[style=plain]
def __init__(self, kind, chars=0, lines=0)
def Step(self, name, chars=None)
def Record(self, name, wall, chars)
def Add(self, name, seconds, chars=0)
def Carry(self, name, seconds, chars=0)
def Total(self)
def AsDict(self)
def AsJson(self)
def Report(self)
\end{lstlisting}
\Q{Method}{Current implementation}
Measure request wall time and retain earlier generation work when drawing is requested later. The excerpt below is from the current definition; the complete function is in the source snapshot. Its constants and branches, rather than an older signature, govern this path.


\begin{lstlisting}[style=plain]
class Run:
    """See source snapshot for the original docstring."""

    def __init__(self, kind, chars=0, lines=0):
        self.kind = kind
        self.chars = chars
        self.lines = lines
        self.stages = []
        self.external = 0.0
        self.t0 = time.perf_counter()

    def Step(self, name, chars=None):
        return Stage(self, name, self.chars if chars is None else chars)

    def Record(self, name, wall, chars):
        self.stages.append(dict(stage=name, seconds=round(wall, 4),
                                characters=chars,
                                secondsPerChar=round(wall/chars, 5) if chars else None))

    def Add(self, name, seconds, chars=0):
        self.Record(name, float(seconds), chars)

    def Carry(self, name, seconds, chars=0):
        self.external += float(seconds)
        self.Record(name, float(seconds), chars)

    def Total(self):
        return (time.perf_counter()-self.t0)+self.external

    def AsDict(self):
        total = self.Total()
        accounted = sum(s["seconds"] for s in self.stages)
        out = dict(
            kind=self.kind,
            characters=self.chars,
            lines=self.lines,
            totalSeconds=round(total, 3),
            measuredSeconds=round(accounted, 3),
            overheadSeconds=round(max(0.0, total-accounted), 3),
            secondsPerCharacter=round(total/self.chars, 4) if self.chars else None,
            secondsPerLine=round(total/self.lines, 3) if self.lines else None,
            charactersPerSecond=round(self.chars/total, 2) if self.chars and total else None,
            stages=self.stages,
        )
        formula = ALLOWANCE.get(self.kind)
        if formula and self.chars:
            limit = formula(self.chars)
            out["limitSeconds"] = round(limit, 2)
            out["limitFormula"] = ALLOWANCE_TEXT.get(self.kind, "")
            out["withinLimit"] = bool(total <= limit)
            out["fractionOfLimit"] = round(total/limit, 3)
            out["secondsSpare"] = round(limit-total, 2)
        return out

    def AsJson(self):
        return json.dumps(self.AsDict(), indent=1)

    def Report(self):
# Excerpt ends; full definition continues in bundled source.
\end{lstlisting}
\Q{Justification}{Interpret the method rather than infer a benchmark}
perf\_counter differences measure elapsed time; Carry adds external work while Add only labels work already inside the wall interval. This avoids double counting stages. \cite{pythontime,code2026}.

